%	현대물리실험 보고서
%	실험5
%	202100973 이승엽

%----------------------------------------------------------------------------------------
%	PACKAGES AND DOCUMENT CONFIGURATIONS
%----------------------------------------------------------------------------------------

\documentclass[a4paper, 10pt, nanum]{CSUniSchoolLabReport}
% use UTF8 encoding
\usepackage[utf8]{inputenc}
% use KoTeX package for Korean
\usepackage{kotex}
% use ams' math font
\usepackage{amsmath}
\usepackage{hyperref}
% use table
\usepackage{graphicx}
\usepackage{multirow}
\usepackage{indentfirst}
\usepackage{setspace}
\usepackage{enumitem}
\usepackage{wrapfig}
\usepackage{epstopdf}
% use captions
\usepackage{caption}
\usepackage{subcaption}
\usepackage{blindtext}
% use layout frame
\usepackage{showframe}

% Setting the graphics path
\graphicspath{{Figures/}}

% only for show page layout
\renewcommand\ShowFrameLinethickness{0.25pt}
\renewcommand*\ShowFrameColor{\color{red}}

\setlength{\parindent}{0.2in} % 들여쓰기 길이 설정
\setlength{\parskip}{3mm} % 문단 간의 간격 조절
\setstretch{1.5} % 줄간격
% \graphicpath{{Figures/}} % fig 경로 설정

\addbibresource{reference.bib} % Bibliography file (located in the same folder as the template)

%----------------------------------------------------------------------------------------
%	REPORT INFORMATION
%----------------------------------------------------------------------------------------

\title{현대물리실험 실험5 보고서 \\ 전자의 회절 및 XRD} % Report title

\author{\textsc{Department of Physics} 202100973 이승엽}

\date{\today}

%----------------------------------------------------------------------------------------

\begin{document}

\maketitle % Insert the title, author and date using the information specified above

\begin{center}
	\begin{tabular}{l r}
		Date Performed: & May 30, 2023 \\ % Date the experiment was performed
		Partners: & 202100969 이규리 \\ % Partner names
		& 202100989 한누리 \\
		Instructor: & Professor 이기주 \\ % Instructor/supervisor
		Typesetting: & LaTeX \\
		Datafitting: & OriginLab \& Python \\
	\end{tabular}
\end{center}

%----------------------------------------------------------------------------------------
%	ABSTRACT
%----------------------------------------------------------------------------------------

\maketitle
% \begin{abstract}
% 	This report ...
% \end{abstract}

%----------------------------------------------------------------------------------------
%	INTRODUCTION
%----------------------------------------------------------------------------------------

\section{Introduction}

\begin{enumerate}[label=\arabic*.]
	\item 전자의 회절 격자에 대해 조사한다.
	\item 회절 무늬를 관측하여 흑연의 interplanar spacing을 측정한다.
	\item X-선을 사용하여 브래그 각을 측정하고, 그 특성 동 선의 에너지 값을 계산하고, 비교한다.
\end{enumerate}



%----------------------------------------------------------------------------------------
%	THEORY
%----------------------------------------------------------------------------------------

\section{Theory}

	파동으로서의 전자 : 운동량 $P$를 갖고 있는 전자의 파장은 드브로이 방정식에 의하면
	\begin{align}
		\lambda = \frac{h}{P}
	\end{align}
	이다. 또 운동량 p는 전자를 가속한 전압이 V라 할 때 운동 에너지의 식에서부터 구할 수 있다.
	\begin{align}
		K = \frac{P^2}{20} = eV
	\end{align}
	로 두 수식을 연립하면
	\begin{align}
		\lambda = \frac{h}{\sqrt{2meV}}
	\end{align}
	가 된다. 참고로 여기서 사용되는 전압 범위에서는 상대론적 질량과 고전적인 질량의 차이는 0.5\% 보다 작다.

	회절 무늬: 전자가 구리(금속입자:Cu) grating에 입혀진 다결정 흑연박막에 부딪치면 브레그 조건에 따라 반사된다. $n$이 자연수 1, 2, 3, ...일 때,
	\begin{align}
		n\lambda = 2d \sin\theta
	\end{align}
	가 된다.

	흑연 결정의 구조와 회절격자 사이의 거리는 Fig. 1, 2와 같다
	\begin{figure}[ht!]
		\centering
		\includegraphics[width=5cm]{Figures/diamond-graphite.png}
		\caption{다이아몬드와 흑연 결정(graphite)의 구조 차이.}
		\label{fig:diamond-graphite}
	\end{figure}
	\begin{figure}[ht!]
		\centering
		\includegraphics[width=5cm]{Figures/graphite.png}
		\caption{흑연 격자 간 거리.}
		\label{fig:graphite}
	\end{figure}

	다결정 측연에서는 각각의 층(layer)사시의 결합(bond)이 깨져서 그 방향이 제멋대로이다. 따라서 전자빔(또는 전자선)은 원뿔모양의 cone처럼 퍼져서 진행하며 회절무늬를 형광판에 원으로 나타난다. 브래그각 $\theta$는 원 모양의 회절 무늬의 반지름에서 측정할 수 있다. (단, Fig. 3에서 원의 특성에 의해 $\alpha = 2\theta$임을 알고 있어야 한다.)
	\begin{figure}[ht!]
		\centering
		\includegraphics[width=5cm]{Figures/experiment.png}
		\caption{실험장치의 개략도.}
		\label{fig:experiment}
	\end{figure}
	따라서
	\begin{align}
		\sin 2\alpha = \frac{D_n}{R}
	\end{align}
	이 된다. 여기서 $R$은 유리 전구의 반지름으로 $2R = L$.이다. 만약 $\alpha$가 매우 작다면 삼각함수의 특성에 의해
	\begin{align}
		D_n = \frac{2Rn\lambda}{d}
	\end{align}

	XRD 분석이란 X-ray를 원하는 시편에 회절 시켜 시편의 내부 정보를 그래프로 나타내는 분석법이다. 보통 최종 결과물은 아래의 그림 그래프와 같다. X축은 $2\theta$로 Y축은 intensity라는 값으로 그래프로 나타내며 여러 상들의 고유한 각도에서 peak값이 나타나 시편 속에 어떠한 phase로 이루어져 있는지 확인할 수 있는 분석 방법이다.
	\begin{figure}[ht!]
		\centering
		\includegraphics[width=5cm]{Figures/exam_xrd.png}
		\caption{}
		\label{fig:exam_xrd}
	\end{figure}

%----------------------------------------------------------------------------------------
%	EXPERIMENTAL METHOD
%----------------------------------------------------------------------------------------

\section{Experimental Method}

\subsection*{실험 1. 전자의 회절}

	\begin{enumerate}[label=\arabic*.]
		
		\item F1과 F2를 110V 전용 high-voltage power supply로 10 kV 걸어준다.
		\item C와 X 커넥터를 주황색 power supply의 음극에 연결한다.
		\item A 커넥터를 주황색 power supply의 양극에 연결한다.
		\item 주황색 power supply의 전압을 서서히 올려준다. (이 전압을 accelerating voltage $U$라고 한다.)
		\item Accelerating voltage $U$를 2 kV에서 5 kV까지 0.5 kV간격으로 올려주며 스크린에 비춘 원형고리의 지름 $D_1$과 $D_2$를 버니어 켈리퍼스로 측정한다.
	\end{enumerate}


% \subsection{STM 측정용 팁 만들기}
% 	\begin{enumerate}[label=\arabic*.]
% 		\item Pt/Ir wire의 일부를 잘라 사용한다. 가장 어려운 부분이므로 신중하게 한다.
% 		Magnifying cover를 The NaioSTM system에서 분리한 NaioSTM tool set에서 the wire cutters, the Flat nose pliers 그리고 the pointed tweezers를 에탄올로 닦는다.		
% 		\item the pointed tweezers를 사용해 '오래된 tip'을 제거한다. 만약, '오래된 tip'이 충분히 길다면 그 '오래된 tip'을 잘라 쓸 수도 있다
% 		\item tip으로 쓰게 될 와이어를 꺼낸다.
% 		\item pliers를 왼손에 쥐고 꽉 잡아 흔들리지 않도록 한다.
% 		\item 그 상태에서 cutters로 약 4 mm정도의 길이로 비스듬히 자른다. 주의할 점은 wire를 싹뚝 잘라내는 것이 아니라 깍아 내듯이 잘라야 한다. tip의 끝이 굉장히 날카로워야 하기 때문에 자르는 것이 아니라 찢어내듯이 잘라야 한다. 그렇지 않으면 선명한 신호를 얻어내지 못한다.
% 		\item the pointed tweezers를 이용해 The NaioSTM system 본체 안쪽에 고정시킨다. 팁이 0.2~0.3 cm정도 튀어나오도록 장착한다.
% 		\item Magnifying cover로 닫고 확대경으로 다음과 같이 장착되었는지 팁 끝부분을 확인한다.
% 		\item 측정샘플을 홀더에 고정하기 위해 NaioSTM tool set에서 홀더를 꺼낸다.
% 		\item 손에 라텍스 장갑을 써서 샘플이 오염되지 않도록 한다.
% 		\item NaioSTM tool set에서 트위져와 측정샘플을 꺼내어 홀더에 장착한다. (홀더에 자성으로 고정이 되므로 끝 부분을 잡고 살짝 놓는다.  장갑을 껴도, 샘플 holder의 금속 부분은 만지지 않도록 해야 합니다.)
% 		\item 측정샘플을 장착한 홀더를 The NaioSTM system헤드에 끼운다.
% 		\begin{figure}[htb!]
% 			\centering
% 			\begin{minipage}{.5\textwidth}
% 				\centering
% 				\includegraphics[width=5cm]{Figures/STM_TIP1.jpg}
% 				\captionof{figure}{Tip mounted on the STM.}
% 				\label{fig:STM_TIP1}
% 			\end{minipage}%
% 			\begin{minipage}{.5\textwidth}
% 				\centering
% 				\includegraphics[width=5cm]{Figures/STM_TIP2.jpg}
% 				\captionof{figure}{Tip mounted on the STM.}
% 				\label{fig:STM_TIP2}
% 			\end{minipage}
% 		\end{figure}
% 		\item 뚜껑을 닫고 팁과 샘플과의 거리가 0.1 cm이내인지 확인한다.
% 		\item Magnifying cover를 덮고 the flat pliers와 기타 도구들을 정리한다.
% 	\end{enumerate}

%----------------------------------------------------------------------------------------
%	RESULT AND DISCUSSION
%----------------------------------------------------------------------------------------

\section{Result and Discussion}

	% STM 측정 프로그램 셋팅을 다음과 같이 설정한다.
	
	% \begin{table*}[htb!]
	% 	\label{tab:2}
	% 	\centering
	% 	\resizebox{\textwidth}{!}{
	% 	\begin{tabular*}{\textwidth}{@{\extracolsep{\fill}}ccc|ccc}
	% 			\noalign{\smallskip}\noalign{\smallskip}\hline\hline
	% 		& Size (nm) & & Point/Line & Time (s) \\
	% 		\hline
	% 		& 200 & & 128 & 0.2 \\
	% 		& 2 & & 512 & 0.5 \\
	% 		\hline
	% 		\hline
	% 	\end{tabular*}
	% 	}
	% \end{table*}

	% Tip을 깎아서 장착 후, Sample HOPG를 장착한다. Sample의 적당한 위치에 맞춰 tip 위치를 조정한다. 프로그램을 사용하여 sample의 위치를 미세 조정한다.

	% Advanced/Retract, approach/withdraw를 사용하여 Tunneling하는 순간까지 이동시킨다. Sample size에 맞춰 측정한다.

	% \begin{figure}[htb!]
	% 	\centering
	% 	\includegraphics[width=5cm]{Figures/Scan_200nm.jpg}
	% 	\caption{Scanned at 200 nm with HOPG.}
	% 	\label{fig:Scan_200nm}
	% \end{figure}

	% 200 nm로 측정했을 때, 미끄러지는 노이즈가 생겼다. 이렇게 될 경우, sample size를 줄이면 노이즈가 급격히 심해질 수 있다. (Figures~4)

	% \begin{figure}[htb!]
	% 	\centering
	% 	\begin{minipage}{.5\textwidth}
	% 		\centering
	% 		\includegraphics[width=5cm]{Figures/Scan_2nm_1.jpg}
	% 		\label{fig:Scan_2nm1}
	% 	\end{minipage}%
	% 	\begin{minipage}{.5\textwidth}
	% 		\centering
	% 		\includegraphics[width=5cm]{Figures/Scan_2nm_2.jpg}
	% 		\label{fig:Scan_2nm2}
	% 	\end{minipage}
	% 	\caption{Scanned at 2 nm with HOPG.}
	% \end{figure}

	% 2 nm로 측정했을 때, 분자 사이즈의 10배 정도의 size로 측정되기에 한 line에 10개 정도의 분자가 보일 것이다. HOPG는 Graphite이기에 육각형 튜브 구조이다. 측정이 잘 되었다면 육각형의 원자 구조를 볼 수 있었겠으나, tip이 잘 만들어지지 않아 노이즈가 생겨 Figures~5와 같이 육각 구조가 이어져 있는 line으로 보인다. (Figures~5)


	% \begin{figure}[htb!]
	% 	\centering
	% 	\begin{minipage}{.5\textwidth}
	% 		\centering
	% 		\includegraphics[width=5cm]{Figures/Scan_2nm_3D1.jpg}
	% 		\label{fig:Scan_2nm_3D1}
	% 	\end{minipage}
	% 	\begin{minipage}{.5\textwidth}
	% 		\centering
	% 		\includegraphics[width=5cm]{Figures/Scan_2nm_3D2.jpg}
	% 		\label{fig:Scan_2nm_3D2}
	% 	\end{minipage}
	% 	\caption{3D scanned at 2nm with HOPG.}
	% \end{figure}

	% 3D로 봤을 때는 2D linear 형태보다 잘 보임을 알 수 있다. Tunneling이 잘 되는 부분이 매우 잘 보인다. 다만 Figures~5와 같이 육각 구조가 보이지는 않았다. (Figures~6)

%----------------------------------------------------------------------------------------
%	CONCLUSIONS
%----------------------------------------------------------------------------------------

\section{Conclusions}

	% STM을 사용하여 HOPG를 Scan하였을 때 대략적인 분자 구조를 확인할 수 있었으며, 3D model로 확인할 경우 tunneling을 볼 수 있었다. Tip을 더 잘 만들고, sample이 깨끗했다면 2 nm or below에서도 노이즈가 적게 측정되어 분자 구조가 보였을 것이다.

%----------------------------------------------------------------------------------------
%	FITTING CODE FOR PYTHON
%----------------------------------------------------------------------------------------

% \section{Fitting code for Python}

%----------------------------------------------------------------------------------------
%	BIBLIOGRAPHY
%----------------------------------------------------------------------------------------

% \printbibliography % Output the bibliography

%----------------------------------------------------------------------------------------

\end{document}